% Generated by tools/build_v6.py; edit the assembly rules, not this file.
\section{Finite certificates for the unrestricted experiment}
\label{sec:genie}

The universal cap in Corollary~\ref{cor:dual-finite-characterization} is a
sufficient theoretical cap. This section gives a different computational
route: upper and lower programs at a chosen cap, with an a posteriori
stopping criterion. The lower program strengthens the one-unit experiment
by revealing its latent magnitude and sign if it reaches that cap.
This reveal-at-cap device has a direct precedent in Lee--Valiant
\cite[Section 5.2]{leevaliant}; the contribution here is its convergent robust
dual bracket and exact stopping certificate. The extra information is
used only in a lower-bound argument. The upper
estimator continues to observe ordinary Bernoulli samples alone.

\subsection{The infinite-cap dual and its compactness argument}

Define
\[
 P_\infty(B)=\inf_{N\ge1}P_N(B),\qquad
 K_\infty(q_L,q_H;q_v,q_w)=\sup_{N\ge1}K_N(q_L,q_H;q_v,q_w).
\]
The latter definition uses only deterministically capped local policies.
No assertion about a general unbounded policy is needed for the following
interchange.

\begin{lemma}[The exact dual persists in the limit]\label{lem:infinite-dual}
For every law and every $B>0$,
\begin{equation}\label{eq:infinite-dual}
 P_\infty(B)^2
 =\max_{q_L,q_H,q_v,q_w\in\mathcal Q_\rho}
 \frac{[\theta(q_H)-\theta(q_L)]_+^2}{4+B K_\infty(q_L,q_H;q_v,q_w)}.
\end{equation}
The maximum is attained and is positive. If $D_\infty$ is the corresponding
midpoint-reference modulus, then
\begin{equation}\label{eq:infinite-comparison}
 D_\infty\le P_\infty\le3D_\infty,
 \qquad R(B)\le P_\infty(B)\le8R(B).
\end{equation}
\end{lemma}

\begin{proof}
The space $\mathcal Q_\rho^4$ is weak-star compact. For each fixed cap,
$K_N$ is lower semicontinuous: every fixed policy has a lower semicontinuous
closed-perspective numerator and a positive continuous work denominator,
and taking their supremum preserves lower semicontinuity. Its dual ratio
$f_N=(\Delta\theta)_+^2/(4+BK_N)$ is upper semicontinuous, with value zero
at an infinite denominator. Indeed, the numerator is continuous and
bounded, and lower semicontinuity forces a diverging denominator when
approaching a point with $K_N=\infty$. The functions $f_N$ decrease with $N$.
If $M=\inf_N\max f_N$, the nonempty compact sets $\{f_N\ge M\}$ are
nested. Their intersection gives a quadruple with $\inf_Nf_N\ge M$.
The reverse inequality is immediate. This proves \eqref{eq:infinite-dual}
and attainment. Positivity follows from \eqref{eq:dual-positive-value}.

The equal-third reference absorption in
Theorem~\ref{thm:midpoint-modulus} holds simultaneously at every cap,
hence also for $K_\infty$. It proves the first two comparisons.
For a fixed midpoint pair, the hybrid Hellinger bound holds at every
finite population cap, with the same coefficient $K_\infty$. Increasing
capped prefixes extends the affinity bound to the full terminal
experiment of any finite-expected-work algorithm. The testing and
interpolation argument of Theorem~\ref{thm:dual-adaptive} therefore gives
$R\ge3D_\infty/8$. The constructive inequality $R\le P_N$ holds for
every $N$. These facts prove \eqref{eq:infinite-comparison}.
\end{proof}

\subsection{Revealing the latent state at the cap}

Fix four priors and write $\delta=q_H-q_L$. Continue to use the prefix
coefficients from Proposition~\ref{prop:dual-bellman}. For a prefix
$s=(n,k)$, let
\[
 \ell_s^\pm(d)=(1/2\pm d)^k(1/2\mp d)^{n-k}.
\]
The ordinary stopping reward is
$r_s=\mathfrak r(p_H(s)-p_L(s),p_v(s))$. Define the latent-reveal reward
\begin{equation}\label{eq:genie-terminal}
 g_s=\int
 \left\{\frac{\delta(d)^2\ell_s^+(d)}{q_v(d)}+
             \frac{\delta(d)^2\ell_s^-(d)}{1-q_v(d)}\right\}\,\mu(dd).
\end{equation}
Terms with zero numerator and denominator are zero; a nonzero numerator
over zero has value infinity. In normalized count coordinates, multiply
$r_s,g_s,p_w(s)$ by $\binom nk$ and use the corresponding Pascal factors.

Let $K_N^+$ be the optimal information/work ratio in the experiment which
uses ordinary observations until it stops before depth $N$, but receives
$(D,S)$ at no extra cost if it reaches depth $N$. Its Bellman recursion is
exactly \eqref{eq:dual-bellman}, with terminal value $g_{Nk}$ in place of
$r_{Nk}$. The root condition is unchanged. Thus $\kappa\ge K_N^+$ exactly
when this modified recursion has root surplus at most zero.

For a finite latent divergence, set
\begin{align}
 C_v&=\int\frac{\delta(d)^2}{q_v(d)(1-q_v(d))}\,\mu(dd),\notag\\
 e_N^v&=C_v-
 \sum_{k=0}^N\binom Nk
 \mathfrak r\bigl(p_H(N,k)-p_L(N,k),p_v(N,k)\bigr).
 \label{eq:genie-residual}
\end{align}

\begin{theorem}[Monotone stopping-ratio brackets]\label{thm:genie-ratio}
For every admissible quadruple of priors, with information ratios
interpreted in the extended nonnegative reals,
\begin{equation}\label{eq:genie-bracket}
 K_N\uparrow K_\infty\le K_N^+,
 \qquad K_{N+1}^+\le K_N^+.
\end{equation}
If, in addition, $C_v<\infty$, then
\begin{equation}\label{eq:genie-error}
 0\le K_N^+-K_N\le e_N^v,\qquad e_N^v\longrightarrow0.
\end{equation}
Consequently $K_N^+\downarrow K_\infty$ under this finite-divergence
hypothesis. Without it, the monotone brackets remain valid, but
pointwise convergence of the upper bracket need not hold.
\end{theorem}

\begin{proof}
Under the variance-reference latent law, define the score
\[
 Z(D,S)=
 \begin{cases}
  \delta(D)/q_v(D),&S=1,\\
  -\delta(D)/(1-q_v(D)),&S=0.
 \end{cases}
\]
If $C_v<\infty$, it is square integrable, has mean zero, and has second
moment $C_v$. The score of any observed transcript is its conditional
expectation. Thus its signed-likelihood information is
$\E_v[\E_v(Z\mid\text{transcript})^2]$.

Take any capped ordinary policy with cap larger than $N$ and stop its paid
observations at $T\wedge N$. On reaching $N$, the genie reveals $(D,S)$;
these values allow simulation of every future Bernoulli observation and
therefore of the original terminal transcript. Data processing cannot
decrease information when passing to the stronger genie transcript, while
the work under $q_w$ can only decrease. This proves $K_\infty\le K_N^+$
by taking the supremum over finite ordinary caps. The same simulation
shows that an earlier reveal is stronger than a later reveal, so $K_N^+$
is nonincreasing. The other monotonicity and $K_N\le K_\infty$ are
immediate. These arguments also apply to infinite divergences using
the extended-value definition.

For a cap-$N$ genie policy, erase the reveal. The resulting ordinary policy
has exactly the same paid work. Its lost information is
\[
 \E_v\left[\mathbf1_{\{T=N\}}
                 \Var_v(Z\mid\mathcal F_N)\right]
 \le\E_v\Var_v(Z\mid\mathcal F_N)=e_N^v.
\]
Here $\mathcal F_N$ is the first $N$ observations of the completed
counterfactual stream. An independent policy seed can be included; it
adds no latent information. Since root-active work under $q_w$ is at
least one, dividing by work and taking suprema gives
$K_N^+\le K_N+e_N^v$.

The infinite Bernoulli stream identifies its mean almost surely and hence
identifies $(D,S)$ because $D>0$. The martingale
$\E_v(Z\mid\mathcal F_N)$ therefore converges to $Z$ in $L^2$.
This proves $e_N^v\to0$. Its explicit formula is precisely
\eqref{eq:genie-residual}.
\end{proof}

\begin{remark}[Pointwise convergence can fail at infinite latent divergence]
Let $\mu(dd)=2\,dd$ on $(0,1/2]$, $\rho=1/4$, and take
$q_L=0$, $q_H=q_w=1/4$, and $q_v(d)=1/2-d$. All four priors are
admissible. Here
\[
 C_v=\frac18\int_0^{1/2}\frac{dd}{1/4-d^2}=\infty.
\]
The policy which always reaches depth $N$ and accepts the latent reveal
has infinite information and work $N$, so $K_N^+=\infty$ for every
finite $N$.

Nevertheless $K_\infty\le3/4$. To verify this for every capped policy,
parameterize the latent state by its Bernoulli mean $p$. The variance
reference has density $2(1-p)$ on $(0,1)$, the work reference has density
$3/2$ below $1/2$ and $1/2$ above $1/2$, and the signed endpoint
difference has density $-1/2$ below $1/2$ and $1/2$ above $1/2$.
For an ordered prefix $s$ of length $n\ge1$ with $k$ heads, put
$u_s=\int_0^1p^k(1-p)^{n-k}\,dp$. Then
\[
 |p_H(s)-p_L(s)|\le u_s/2,\qquad
 p_v(s)=\frac{2(n-k+1)}{n+2}u_s,
\]
and hence
\[
 r_s\le\frac{n+2}{8(n-k+1)}u_s\le\frac{3n}{8}u_s.
\]
Sum over stopped prefixes and the independent seed. If $U$ denotes the
auxiliary uniform law of $p$, the information of any capped policy is
at most $(3/8)\E_U T$. The work-reference density dominates half the
uniform density, so this is at most $(3/4)\E_{q_w}T$. Taking suprema
gives the claimed bound. The auxiliary law $U$ need not be admissible;
it is used only for this domination calculation.

This example does not contradict the optimized convergence in the next
theorem, whose proof first regularizes the variance reference.
\end{remark}

\subsection{A convergent lower program and a stopping rule}

Define the strengthened-experiment dual value
\begin{equation}\label{eq:genie-program}
 G_N(B)^2=
 \sup_{q_L,q_H,q_v,q_w\in\mathcal Q_\rho}
 \frac{[\theta(q_H)-\theta(q_L)]_+^2}{4+B K_N^+(q_L,q_H;q_v,q_w)}.
\end{equation}
As before, a ratio with infinite denominator is zero.

\begin{theorem}[A posteriori minimax brackets]\label{thm:genie-program}
For every law, rarity and budget in the model,
\begin{equation}\label{eq:genie-program-convergence}
 0<G_N(B)\uparrow P_\infty(B),\qquad P_N(B)\downarrow P_\infty(B),
\end{equation}
and at every integer $N\ge1$,
\begin{equation}\label{eq:genie-risk-bracket}
 \boxed{\qquad G_N(B)/8\le R_{\mu,\rho}(B)\le P_N(B).\qquad}
\end{equation}
For finite-support laws both programs have finite conic formulations.
The condition $P_N(B)\le2G_N(B)$ eventually holds and certifies
$R(B)\le P_N(B)\le16R(B)$.
\end{theorem}

\begin{proof}
The preceding theorem gives $G_N\le P_\infty$ and monotonicity. The
constant priors $q_L=0,q_H=q_v=\rho$ have finite latent divergence and
give a strictly positive lower bound for every $G_N$.

To prove convergence, fix a quadruple with finite $K_\infty$ and positive
target gap. Put $m=(q_L+q_H)/2$ and regularize only the variance reference:
\[
 q_v^\epsilon=(1-\epsilon)q_v+\epsilon m,
 \qquad 0<\epsilon<1.
\]
The new transcript reference dominates $(1-\epsilon)$ times the old one,
so
\[
 K_\infty(q_L,q_H;q_v^\epsilon,q_w)
 \le\frac{K_\infty(q_L,q_H;q_v,q_w)}{1-\epsilon}.
\]
Also $|\delta|\le2m$ and $|\delta|\le2(1-m)$, so the new latent score
has absolute value at most $2/\epsilon$. Its latent divergence is finite.
Theorem~\ref{thm:genie-ratio} now makes its genie ratio converge to its
ordinary infinite-cap ratio. Keep $\epsilon>0$ fixed throughout
this limit in $N$; no convergence of the unregularized
quadruple is used. Hence
\[
 \lim_NG_N(B)^2\ge
 \frac{(\Delta\theta)^2}{4+B K_\infty/(1-\epsilon)}.
\]
Let $\epsilon\downarrow0$ and maximize using
Lemma~\ref{lem:infinite-dual}. The limits are taken in this order;
no convergence uniform in $\epsilon$ is asserted. This proves
$G_N\uparrow P_\infty$.
Equation~\eqref{eq:infinite-comparison} gives the risk bracket.
Since $P_\infty>0$, the ratio $P_N/G_N$ tends to one. The stated stopping
criterion therefore succeeds at a finite cap, with no claimed uniform
rate for that cap.
\end{proof}

For an explicitly given law $\mu=\sum_iw_i\delta_{d_i}$, the genie dual
is a finite SOCP. In the price formulation, use endpoint variables
$u_H=\alpha q_H$, $u_L=\alpha q_L$, work-reference variables
$e=\eta q_w$, and a variance reference $q_v$. Their rarity and box
constraints are linear perspective constraints. For each atom introduce
\[
 (u_{H,i}-u_{L,i})^2\le q_{v,i}r_i^+,
 \qquad
 (u_{H,i}-u_{L,i})^2\le(1-q_{v,i})r_i^-.
\]
At terminal states replace the ordinary quadratic-over-linear stopping
cone by the linear bound
\[
 X_s\ge\frac14\sum_iw_i
           \{\ell_s^+(d_i)r_i^++\ell_s^-(d_i)r_i^-\}.
\]
Keep the ordinary stopping cones at earlier states and the same linear
continuation/root inequalities. Maximize
$\sum_iw_i(u_{H,i}-u_{L,i})-\alpha^2-\eta B$.
This is exactly the squared value in \eqref{eq:genie-program} after
eliminating the price and scale. It uses $O(N^2+m)$ variables, apart from
the coefficient matrices. For a general measure the result remains a
variational statement; a finite-support computational representation is
not silently assumed for an arbitrary law.

\subsection{Certificates that can be checked independently}

In practice a feasible primal policy gives an upper bound $U\ge R$, and
a feasible genie dual gives a square-root lower bound $L\le G_N$.
Consequently the \emph{verified} inequality $U\le2L$ is sufficient for
a factor-16 risk certificate. An approximate solver status or an
uncertified objective is not sufficient. Section~\ref{sec:computation}
uses exact rational checks for the finite rational examples.

A selected dual quadruple can also yield a stronger direct risk lower
bound than $L/8$. Absorb its two references into its endpoints with equal
thirds, and let $t$ be the new target separation. If a Bellman certificate
gives $\kappa\ge K_N^+$ for the original quadruple, the absorbed midpoint
pair has $K_\infty\le\kappa$. With $x=B\kappa$, affinity testing gives
\begin{equation}\label{eq:direct-genie-risk}
 R(B)\ge\max_{0\le s\le1}
          \frac{st}{2}\left(1-\frac{s^2x}{4}\right)_+.
\end{equation}
The square of this bound is the rational expression
\begin{equation}\label{eq:direct-genie-risk-squared}
 \begin{cases}
 \displaystyle\frac{t^2}{4}(1-x/4)^2,&0\le x\le4/3,\\[4pt]
 \displaystyle\frac{4t^2}{27x},&x\ge4/3.
 \end{cases}
\end{equation}
This supplies an exact squared lower certificate when the inputs and
the certified price are rational.
